\documentclass{article}

% Language setting
% Replace `english' with e.g. `spanish' to change the document language
\usepackage[english]{babel}

% Set page size and margins
% Replace `letterpaper' with`a4paper' for UK/EU standard size
\usepackage[letterpaper,top=2cm,bottom=2cm,left=3cm,right=3cm,marginparwidth=1.75cm]{geometry}

% Useful packages
\usepackage{amsmath}
\usepackage{graphicx}
\usepackage[colorlinks=true, allcolors=blue]{hyperref}

\title{Title of your Algorithm/Software}
\author{Your Roll numbers}

\begin{document}
\maketitle

\begin{abstract}
Your abstract summarizes your Algorithm/Software/Tool..
\end{abstract}

\section{Main task/Tool purpose}

Briefly explain what question is your tool design for.

\section{Input/Output}

Explain what inputs are expected/allowed and what will be returned by your tool.

\section{Key algorithms}

Details of what was developed.

\subsection{Parameters}

Which are the parameters used? How to choose them? What are the defaults?

\section{Technical details}

What are the dependencies? How to install the software? How to run the software? What are the expected running times? What are the measures of accuracy? 

\section{Division of work}

This section should briefly describe what was the contribution of each member of the team. Use the roll numbers.

\section{References}

Add any essential reference (if needed). Any (consistent) format is fine.

The total length of this report cannot exceed 3 pages.

[This report can have different sections and subsections as appropriate for the project]

\end{document}